% ==========================================
% ФАЙЛ: p9.tex
% БЛОК БИЛЕТОВ 41-45 (Функциональные последовательности и ряды)
% ==========================================

% --- Вопрос 41 ---
\examquestion{41}{Функциональные последовательности: поточечная и равномерная сходимость}
{Сформулируйте определения поточечной и равномерной сходимости. Дайте геометрическое обоснование (критерий $\varepsilon$-полосы). Нарисуйте схему зажима графиков внутри трубчатой окрестности и покажите, почему поточечная сходимость не гарантирует это свойство.}

\paragraph{Суть на пальцах и Базовый ликбез}
\begin{itemize}
    \item \textbf{Ликбез:} Раньше мы складывали числа, теперь работаем с графиками. В функциональной последовательности $\{f_n(x)\}$ или ряде $\sum u_n(x)$ каждый член — это функция. Сходимость ряда всегда исследуется через последовательность его частичных сумм $S_n(x) = \sum_{k=1}^n u_k(x)$. То есть вопрос сходимости ряда полностью сводится к поведению последовательности $S_n(x)$.
    \item \textbf{Что от чего зависит (Жесткий ориентир):}
    \begin{itemize}
        \item \textbf{Поточечная ($f_n \to f$):} Номер $N = N(\varepsilon, x)$. Скорость сходимости привязана к конкретной точке. В точке $x_1$ график может приблизиться к пределу уже при $n=5$, а в точке $x_2$ ему понадобится $n=10000$.
        \item \textbf{Равномерная ($f_n \rightrightarrows f$):} Номер $N = N(\varepsilon)$. Полная независимость от $x$. Существует один общий номер $N$, после которого ВСЕ точки графика одновременно падают в целевой коридор.
    \end{itemize}
\end{itemize}

\paragraph{1. Строгие определения (Игра кванторов)}
Пусть последовательность функций $\{f_n(x)\}$ задана на множестве $X$.
\begin{itemize}
    \item \textbf{Поточечная сходимость:}
    $$ \forall \varepsilon > 0 \quad \forall x \in X \quad \exists N(\varepsilon, x) \in \mathbb{N} : \quad \forall n > N \implies |f_n(x) - f(x)| < \varepsilon $$
    
    \item \textbf{Равномерная сходимость:}
    $$ \forall \varepsilon > 0 \quad \exists N(\varepsilon) \in \mathbb{N} \quad \forall x \in X : \quad \forall n > N \implies |f_n(x) - f(x)| < \varepsilon $$
\end{itemize}

\paragraph{2. Геометрическое обоснование и «трубчатая окрестность»}
Неравенство $|f_n(x) - f(x)| < \varepsilon$ задает $\varepsilon$-полосу (коридор шириной $2\varepsilon$) вокруг предельной функции:
$$ f(x) - \varepsilon < f_n(x) < f(x) + \varepsilon $$
* **При $\rightrightarrows$:** Начиная с номера $N(\varepsilon)$, график $f_n(x)$ целиком и без остатка упакован внутрь этой трубы на всем множестве $X$.
* **При поточечной (без $\rightrightarrows$):** Каждая отдельная точка со временем зайдет в трубу, но зафиксировать график целиком внутри нее невозможно — всегда будет оставаться «вредный» кусок (выброс/горб), который вылезает наружу.

\vspace{0.5em}
\begin{center}
\begin{tikzpicture}[scale=1.2, every node/.style={scale=0.9}]
    \draw[->, thick] (-0.5, 0) -- (4, 0) node[right] {$x$};
    \draw[->, thick] (0, -0.5) -- (0, 3) node[above] {$y$};
    
    % Предельная функция f(x)
    \draw[ultra thick, blue, domain=0:3.5, samples=100] plot (\x, {0.5*(\x-1.5)^2 + 1}) node[right] {$f(x)$};
    
    % Трубчатая окрестность
    \draw[dashed, blue!50, domain=0:3.5, samples=100] plot (\x, {0.5*(\x-1.5)^2 + 1 + 0.5});
    \draw[dashed, blue!50, domain=0:3.5, samples=100] plot (\x, {0.5*(\x-1.5)^2 + 1 - 0.5});
    
    % f_n(x) которая уложилась (равномерная)
    \draw[thick, red, domain=0:3.5, samples=100] plot (\x, {0.5*(\x-1.5)^2 + 1 + 0.3*sin(\x*400)}) node[below right] {$f_{n>N}(x)$};
    
    % "Выброс" - нарушение равномерности
    \draw[thick, orange, domain=0:3.5, samples=100] plot (\x, {0.5*(\x-1.5)^2 + 1 + 1.2*exp(-50*(\x-2.5)^2)});
    \node[orange] at (2.5, 2.5) {Выброс (нарушение $\rightrightarrows$)};
    
    \draw[<->] (0.5, 1.5) -- (0.5, 2) node[midway, left] {$\varepsilon$};
\end{tikzpicture}
\end{center}
\vspace{0.5em}
% --- Вопрос 41 ---
\begin{tcolorbox}[colback=gray!5, colframe=black, boxrule=0.5pt, arc=0pt]
\textbf{Резюмируем Вопрос 41 (Поточечная и равномерная сходимость):} \\
\textbf{Тезис:} Поточечная сходимость ($f_n \to f$) привязана к точке ($N = N(\varepsilon, x)$); равномерная ($f_n \rightrightarrows f$) одинакова для всех ($N = N(\varepsilon)$) и зажимает графики в $\varepsilon$-полосу. \\
\textbf{Ход:} Записываем кванторы: для равномерной сходимости квантор $\exists N(\varepsilon)$ стоит перед $\forall x$, что гарантирует выбор единого номера для всего множества $X$. \\
\textbf{Трюк:} Геометрически равномерная сходимость целиком упаковывает график $f_n(x)$ внутрь трубчатой окрестности $f(x) \pm \varepsilon$ при всех $n > N$. \\
\textbf{Финал:} При поточечной сходимости графики могут содержать локальные «выбросы» (горбы), которые вылезают наружу из трубы при любых номерах $n$.
\end{tcolorbox}



% --- Вопрос 42 ---
\examquestion{42}{Исследование равномерной сходимости и супремум-критерий}
{Сформулируйте и докажите теорему о супремум-критерии равномерной сходимости. Разберите пример «бегущего горба» ($f_n(x) = x^n - x^{2n}$ на $[0,1]$), доказав наличие поточечной сходимости при отсутствии равномерной.}

\paragraph{Суть на пальцах}
Чтобы проверить, уложился ли весь график в $\varepsilon$-трубу, не нужно проверять каждую точку. Достаточно найти на графике \textbf{самую дальнюю} от предела точку (максимальное отклонение, супремум) и посмотреть на её высоту $\rho_n$. Если этот "горб" с ростом $n$ сплющивается до нуля ($\lim \rho_n = 0$), то график уложился в трубу. Если же горб бежит по оси, не теряя высоты, равномерной сходимости нет.

\paragraph{1. Формулировка и доказательство супремум-критерия}
Обозначим $\rho_n = \sup_{x \in X} |f_n(x) - f(x)|$. Тогда:
$$ f_n(x) \rightrightarrows f(x) \text{ на } X \iff \lim_{n \to \infty} \rho_n = 0 $$

\begin{proof}
$(\implies)$ Пусть $f_n(x) \rightrightarrows f(x)$. По определению, для $\forall \varepsilon > 0 \,\, \exists N$, что для $n > N$ и $\forall x \in X$ выполнено $|f_n(x) - f(x)| < \varepsilon/2$. Переходя к точной верхней грани по всем $x$, получаем $\sup_{x} |f_n(x) - f(x)| \le \varepsilon/2 < \varepsilon$. То есть $\rho_n < \varepsilon$ для всех $n > N$. Это и означает, что $\lim \rho_n = 0$.

$(\impliedby)$ Пусть $\lim \rho_n = 0$. По определению предела числовой последовательности, $\forall \varepsilon > 0 \,\, \exists N$, что для $n > N$ выполнено $\rho_n < \varepsilon$. Но так как $\rho_n$ — это супремум (то есть верхняя оценка для всех отклонений), то для любого $x \in X$ справедливо: $|f_n(x) - f(x)| \le \rho_n < \varepsilon$. Это в точности определение равномерной сходимости.
\end{proof}

\paragraph{2. Классический контрпример «Бегущий горб»}
Рассмотрим $f_n(x) = x^n - x^{2n}$ на отрезке $X = [0, 1]$.
1. \textbf{Поточечная сходимость:} 
При фиксированном $x \in [0, 1)$ имеем $x^n \to 0$, поэтому $\lim_{n \to \infty} (x^n - x^{2n}) = 0$.
При $x = 1$ имеем $1^n - 1^{2n} = 1 - 1 = 0$.
Предельная функция $f(x) = 0$ для всех $x \in [0, 1]$. Поточечная сходимость есть.

2. \textbf{Нарушение равномерной сходимости (через $\sup$):}
Найдем максимум отклонения $|f_n(x) - 0|$. Возьмем производную:
$$ f'_n(x) = n x^{n-1} - 2n x^{2n-1} = n x^{n-1} (1 - 2x^n) $$
Приравниваем к нулю: $x^n = \frac{1}{2}$, откуда точка максимума $x_0 = \left(\frac{1}{2}\right)^{1/n}$.
Вычисляем высоту горба в этой точке:
$$ \rho_n = f_n(x_0) = \left(\frac{1}{2}\right) - \left(\frac{1}{2}\right)^2 = \frac{1}{2} - \frac{1}{4} = \frac{1}{4} $$
Перейдем к пределу: $\lim_{n \to \infty} \rho_n = \lim_{n \to \infty} \frac{1}{4} = \frac{1}{4} \ne 0$.
По супремум-критерию равномерной сходимости \textbf{нет}. Горб сплющивается по ширине и бежит к $x=1$, но его абсолютная высота всегда остается равной $1/4$, пробивая любую $\varepsilon$-трубу, если $\varepsilon < 1/4$.
% --- Вопрос 42 ---
\begin{tcolorbox}[colback=gray!5, colframe=black, boxrule=0.5pt, arc=0pt]
\textbf{Резюмируем Вопрос 42 (Супремум-критерий и бегущий горб):} \\
\textbf{Тезис:} Равномерная сходимость на $X$ равносильна тому, что максимальное отклонение $\rho_n = \sup_{X} |f_n(x) - f(x)|$ стремится к нулю при $n \to \infty$. \\
\textbf{Ход:} Для исследования «бегущего горба» $f_n(x) = x^n - x^{2n}$ на $[0,1]$ находим поточечный предел ($f(x)=0$) и ищем глобальный максимум через производную. \\
\textbf{Трюк:} Точка максимума смещается к единице $x_0 = (1/2)^{1/n}$, но высота горба остается жесткой константой: $\rho_n = f_n(x_0) = 1/4$. \\
\textbf{Финал:} Так как $\lim \rho_n = 1/4 \ne 0$, супремум-критерий провален $\implies$ равномерной сходимости нет, хотя поточечная сходимость идеальна.
\end{tcolorbox}



% --- Вопрос 43 ---
\examquestion{43}{Критерий Коши для функциональных последовательностей}
{Сформулируйте и проведите доказательство критерия Коши равномерной сходимости функциональной последовательности. На этапе необходимости обоснуйте переход через неравенство треугольника, а на этапе достаточности постройте предельную функцию $f(x)$ и докажите равномерность стремления к ней.}

\paragraph{1. Формулировка теоремы}
Последовательность $f_n(x)$ сходится равномерно на $X$ $\iff$ она равномерно фундаментальна на $X$:
$$ \forall \varepsilon > 0 \quad \exists N \in \mathbb{N} \quad \forall x \in X : \quad \forall n, m > N \implies |f_n(x) - f_m(x)| < \varepsilon $$

\paragraph{2. Доказательство}
\begin{proof}
\textbf{Необходимость ($\implies$):} 
Пусть $f_n \rightrightarrows f$ на $X$. Зафиксируем $\varepsilon > 0$. По определению равномерной сходимости найдется номер $N$, что для всех $k > N$ и всех $x \in X$ выполнено $|f_k(x) - f(x)| < \frac{\varepsilon}{2}$. 
Возьмем любые $n, m > N$. Применим классическое неравенство треугольника:
$$ |f_n(x) - f_m(x)| = |f_n(x) - f(x) + f(x) - f_m(x)| \le |f_n(x) - f(x)| + |f_m(x) - f(x)| < \frac{\varepsilon}{2} + \frac{\varepsilon}{2} = \varepsilon $$
Необходимость доказана.

\textbf{Достаточность ($\impliedby$):}
Пусть выполнено условие Коши.
\textit{Шаг А: Построение функции $f(x)$.} Зафиксируем любую точку $x_0 \in X$. Тогда числовая последовательность $y_n = f_n(x_0)$ является фундаментальной в $\mathbb{R}$. В силу полноты вещественной прямой, любая фундаментальная последовательность имеет конечный предел. Обозначим это предельное число как $f(x_0)$. Проделав это для каждой точки $x \in X$, мы поточечно построим предельную функцию $f(x)$.

\textit{Шаг Б: Доказательство равномерности.} В условии Коши зафиксируем $n > N$, а индекс $m$ устремим в бесконечность ($m \to \infty$).
Так как при фиксированном $x$ последовательность $f_m(x) \to f(x)$, то, переходя к пределу в нестрогом неравенстве, получаем:
$$ \lim_{m \to \infty} |f_n(x) - f_m(x)| \le \varepsilon \implies |f_n(x) - f(x)| \le \varepsilon $$
Ключевой момент: номер $N$ был выбран в условии Коши и зависел \textbf{только} от $\varepsilon$, но не от $x$. Полученное неравенство работает сразу для всех $x \in X$. Это и есть определение равномерной сходимости. Достаточность доказана.
\end{proof}

% --- Вопрос 43 ---
\begin{tcolorbox}[colback=gray!5, colframe=black, boxrule=0.5pt, arc=0pt]
\textbf{Резюмируем Вопрос 43 (Критерий Коши):} \\
\textbf{Тезис:} Последовательность сходится равномерно $\iff$ она равномерно фундаментальна: $|f_n(x) - f_m(x)| < \varepsilon$ для всех $n,m > N$ сразу на всем множестве $X$. \\
\textbf{Ход:} В необходимости применяем неравенство треугольника, искусственно добавляя предел: $|f_n - f_m| \le |f_n - f| + |f_m - f| < \varepsilon/2 + \varepsilon/2 = \varepsilon$. \\
\textbf{Трюк:} В достаточности сначала поточечно строим предельную функцию $f(x)$ за счет полноты $\mathbb{R}$, а затем в условии Коши устремляем индекс $m \to \infty$. \\
\textbf{Финал:} Нестрогий предельный переход дает $|f_n(x) - f(x)| \le \varepsilon$, где номер $N$ зависит только от $\varepsilon$, что гарантирует равномерность стремления.
\end{tcolorbox}



% --- Вопрос 44 ---
\examquestion{44}{Мажорантный признак Вейерштрасса}
{Сформулируйте мажорантный признак Вейерштрасса для рядов. Проведите его доказательство (через критерий Коши и неравенство треугольника). Объясните, почему этот признак достаточный, но не необходимый.}

\paragraph{Суть на пальцах}
Представь, у тебя есть функциональный ряд $\sum u_n(x)$. Если ты можешь накрыть каждый его член сверху «числовой крышей» (мажорантой), которая не зависит от $x$ ($|u_n(x)| \le a_n$), и сумма этих крыш $\sum a_n$ сходится, то твой функциональный ряд спрятан надежно. Признак Вейерштрасса говорит, что в таком случае ряд сходится \textbf{абсолютно и равномерно}.

\paragraph{Зачем здесь Критерий Коши?}
Мы не знаем итоговую сумму ряда $S(x)$. Критерий Коши позволяет доказать сходимость, работая только с хвостами ряда. Он говорит: если любой кусок ряда в бесконечности (от $n$ до $n+p$) меньше $\varepsilon$ \textbf{сразу для всех $x$}, то ряд сходится равномерно. Мажоранта $a_n$ как раз отвязывает этот хвост от $x$.

\paragraph{1. Формулировка и доказательство}
Пусть для функционального ряда $\sum_{n=1}^{\infty} u_n(x)$ на множестве $X$ существует сходящийся числовой ряд с положительными членами $\sum_{n=1}^{\infty} a_n$, такой что $\forall x \in X$ выполняется неравенство $|u_n(x)| \le a_n$. Тогда ряд $\sum u_n(x)$ сходится абсолютно и равномерно.

\begin{proof}
1. \textbf{Числовой хвост:} Так как $\sum a_n$ сходится, по критерию Коши для числовых рядов:
$$ \forall \varepsilon > 0 \quad \exists N \in \mathbb{N} : \quad \forall n > N, \ \forall p \in \mathbb{N} \implies \sum_{k=n+1}^{n+p} a_k < \varepsilon $$

2. \textbf{Функциональный хвост:} Рассмотрим кусок функционального ряда от $n+1$ до $n+p$. Навесим модуль и применим неравенство треугольника:
$$ \left| \sum_{k=n+1}^{n+p} u_k(x) \right| \le \sum_{k=n+1}^{n+p} |u_k(x)| $$

3. \textbf{Замена на крышу:} Заменяем модули $|u_k(x)|$ на мажоранту $a_k$:
$$ \sum_{k=n+1}^{n+p} |u_k(x)| \le \sum_{k=n+1}^{n+p} a_k < \varepsilon $$

\textbf{Итог:} Мы получили $\left| \sum_{k=n+1}^{n+p} u_k(x) \right| < \varepsilon$. Это неравенство выполнено для всех $x \in X$ одновременно (так как $N$ зависело только от $a_k$). По критерию Коши для функциональных рядов это доказывает равномерную сходимость. Наличие модуля в доказательстве доказывает и абсолютную сходимость.
\end{proof}

\paragraph{2. Достаточный, но не необходимый}
\begin{itemize}
    \item \textbf{Достаточный:} Если мажоранта $a_n$ есть, равномерная сходимость гарантирована.
    \item \textbf{Не необходимый:} Если мажоранты нет, равномерная сходимость всё равно может быть. 
\end{itemize}
\textbf{Контрпример:} Условно сходящийся ряд $\sum_{n=1}^{\infty} \frac{(-1)^n}{n+x^2}$. По признаку Дирихле он сходится равномерно на $\mathbb{R}$. Но признак Вейерштрасса к нему неприменим: если взять модуль, получим $\frac{1}{n+x^2}$, что при $x=0$ дает $\frac{1}{n}$ (расходящийся гармонический ряд). Мажоранты не существует, а равномерная сходимость есть.

% --- Вопрос 44 ---
\begin{tcolorbox}[colback=gray!5, colframe=black, boxrule=0.5pt, arc=0pt]
\textbf{Резюмируем Вопрос 44 (Признак Вейерштрасса):} \\
\textbf{Тезис:} Если члены ряда ограничены числовой мажорантой ($|u_n(x)| \le a_n$) и ряд $\sum a_n$ сходится, то функциональный ряд сходится абсолютно и равномерно. \\
\textbf{Ход:} Оцениваем кусок функционального ряда через неравенство треугольника: $|\sum_{k=n+1}^{n+p} u_k(x)| \le \sum_{k=n+1}^{n+p} |u_k(x)|$. \\
\textbf{Трюк:} Заменяем функциональный «хвост» на числовой: $\sum |u_k(x)| \le \sum_{k=n+1}^{n+p} a_k < \varepsilon$ по критерию Коши для сходящегося числового ряда. \\
\textbf{Финал:} Признак является достаточным, но не необходимым: ряд $\sum \frac{(-1)^n}{n+x^2}$ сходится равномерно по Дирихле, но числовой мажоранты не имеет.
\end{tcolorbox}




% --- Вопрос 45 ---
\examquestion{45}{Свойство непрерывности суммы ряда (Метод трех $\varepsilon$)}
{Сформулируйте и докажите теорему о непрерывности суммы $S(x) = \sum u_n(x)$. Проведите доказательство методом расщепления погрешности на три составляющие («метод трех $\varepsilon$»).}

\paragraph{1. База: определение непрерывности}
Функция $f(x)$ непрерывна в точке $x_0$, если:
$$ \forall \varepsilon > 0 \quad \exists \delta > 0 : \quad \forall x \in X, \ |x - x_0| < \delta \implies |f(x) - f(x_0)| < \varepsilon $$
\textit{Задача:} доказать, что если все члены $u_n(x)$ непрерывны и ряд $\sum u_n(x) \rightrightarrows S(x)$, то его бесконечная сумма $S(x)$ тоже удовлетворяет этому строгому критерию.

\paragraph{2. Доказательство (Метод трёх $\varepsilon$)}
\begin{proof}
Зафиксируем произвольное $\varepsilon > 0$ и точку $x_0$. Нам нужно найти такое $\delta > 0$, чтобы для всех $x$ из $\delta$-окрестности выполнялось $|S(x) - S(x_0)| < \varepsilon$.

\textbf{Шаг 1: Равномерная сходимость (зажимаем хвосты).} 
Так как ряд сходится равномерно, по определению для значения $\frac{\varepsilon}{3}$ найдется ОДИН общий номер $N$, такой что для \textbf{всех} $x$ сразу остаток ряда $R_N(x)$ мал. Значит, мы можем гарантировать малость отклонения истинной суммы $S$ от частичной суммы $S_N$:
$$ |S(x) - S_N(x)| < \frac{\varepsilon}{3} \quad \text{и} \quad |S_N(x_0) - S(x_0)| < \frac{\varepsilon}{3} $$

\textbf{Шаг 2: Непрерывность частичной суммы (выбираем $\delta$).} 
Функция $S_N(x) = \sum_{n=1}^N u_n(x)$ — это конечная сумма непрерывных функций, следовательно, она непрерывна сама по себе. По определению непрерывности, для того же значения $\frac{\varepsilon}{3}$ найдется такое $\delta > 0,$, что при $|x - x_0| < \delta$:
$$ |S_N(x) - S_N(x_0)| < \frac{\varepsilon}{3} $$

\textbf{Шаг 3: Объединение (три эпсилон).} 
Используем неравенство треугольника для полной разности, искусственно добавив и вычтя промежуточные точки $S_N(x)$ и $S_N(x_0)$:
$$ |S(x) - S(x_0)| = |(S(x) - S_N(x)) + (S_N(x) - S_N(x_0)) + (S_N(x_0) - S(x_0))| $$
$$ \le \underbrace{|S(x) - S_N(x)|}_{1\text{-й } \varepsilon/3} + \underbrace{|S_N(x) - S_N(x_0)|}_{2\text{-й } \varepsilon/3} + \underbrace{|S_N(x_0) - S(x_0)|}_{3\text{-й } \varepsilon/3} $$
Итого: при $|x - x_0| < \delta$ получаем $|S(x) - S(x_0)| < \frac{\varepsilon}{3} + \frac{\varepsilon}{3} + \frac{\varepsilon}{3} = \varepsilon$. 
Определение выполнено, сумма ряда $S(x)$ непрерывна в точке $x_0$.
\end{proof}

\paragraph{Анатомия трех $\varepsilon/3$ (Что значит каждая разность)}
\begin{itemize}
    \item $\mathbf{|S(x) - S_N(x)| < \frac{\varepsilon}{3}}$ — \textbf{Вертикальный прыжок от графика к полиному в точке $x$.} Показывает погрешность обрушения бесконечного ряда до конечной суммы в «подвижной» точке. Мал в силу равномерной сходимости.
    \item $\mathbf{|S_N(x) - S_N(x_0)| < \frac{\varepsilon}{3}}$ — \textbf{Горизонтальный шаг вдоль графика полинома $S_N$.} Показывает, насколько сильно меняется конечная частичная сумма при переходе от $x$ к $x_0$. Мал в силу непрерывности конечной суммы (контролируется выбором $\delta$).
    \item $\mathbf{|S_N(x_0) - S(x_0)| < \frac{\varepsilon}{3}}$ — \textbf{Вертикальный прыжок от полинома к графику в точке $x_0$.} Погрешность обрушения ряда в самой фиксированной точке $x_0$. Мал в силу той же равномерной сходимости.
\end{itemize}

\paragraph{Схема «Эстафеты»}
$$ 
|S(x) - S(x_0)| \le \underbrace{|S(x) - S_N(x)|}_{\text{Спуск к частичной сумме } S_N} + \underbrace{|S_N(x) - S_N(x_0)|}_{\text{Переход по непрерывному мосту } S_N} + \underbrace{|S_N(x_0) - S(x_0)|}_{\text{Подъем к истинному значению } S}
$$

\paragraph{Почему важна равномерность и Контрпример} 
Без равномерной сходимости номер $N$ жестко привязан к конкретному иксу: $N = N(\varepsilon, x)$. При приближении $x \to x_0$ скорость сходимости может падать, требуя все больших и больших номеров $N$. Мы потеряем возможность зафиксировать \textbf{один единый} номер для всей окрестности, из-за чего первый хвост будет взрываться, и эстафета разорвется.

\textbf{Классический контрпример (разрушение непрерывности):}
Рассмотрим ряд, частичные суммы которого имеют вид $S_n(x) = x^n$ на отрезке $[0,1]$ (каждый член ряда $u_n(x) = x^n - x^{n-1}$ непрерывен). 
\begin{itemize}
    \item При любом фиксированном $x \in [0, 1)$ предел $S_n(x) = x^n \to 0$.
    \item При $x = 1$ предел $S_n(1) = 1^n \to 1$.
\end{itemize}
Итоговая сумма ряда равна:
$$
S(x) = \begin{cases} 0, & x \in [0, 1) \\ 1, & x = 1 \end{cases}
$$
Все функции внутри ряда были идеально непрерывными, но их бесконечная сумма $S(x)$ получила \textbf{разрыв} в точке $x=1$. Причина — сходимость на $[0,1]$ была поточечной, но не равномерной (графики $x^n$ до последнего «цепляются» за единицу и не могут целиком упасть в горизонтальный коридор вокруг нуля).

% --- Вопрос 45 ---
\begin{tcolorbox}[colback=gray!5, colframe=black, boxrule=0.5pt, arc=0pt]
\textbf{Резюмируем Вопрос 45 (Метод трех $\varepsilon$):} \\
\textbf{Тезис:} Бесконечная сумма $S(x) = \sum u_n(x)$ непрерывна, если все слагаемые непрерывны и ряд сходится равномерно. \\
\textbf{Ход:} Расщепляем погрешность $|S(x) - S(x_0)|$ на три шага: вертикальный спуск к частичной сумме $S_N(x)$, переход по непрерывному мосту к $x_0$, и подъем к $S(x_0)$. \\
\textbf{Трюк:} Равномерная сходимость дает $|S - S_N| < \varepsilon/3$ независимо от $x$, а непрерывность конечной суммы $S_N$ позволяет подобрать $\delta$, чтобы $|S_N(x) - S_N(x_0)| < \varepsilon/3$. \\
\textbf{Финал:} Складывая три составляющие по $\varepsilon/3$, получаем полную непрерывность. Без равномерной сходимости эстафета рвется (пример: $S_n(x) = x^n$ на $[0,1]$).
\end{tcolorbox}